\documentclass[12pt]{article} 

%Paquetes Basicos
\usepackage[spanish]{babel} %identificar idiomas
\usepackage{amsmath} % operaciones matematicas
\usepackage{graphicx} % figuras y graficos

\title{Práctica 1 en LaTeX}
\author{Jesús Emilio Jiménez Aguilar 286666}
\date{\today}

\begin{document}
\maketitle

\section*{Introducción}
\LaTeX Es una herramienta utilizada para la elaboracion de documentos técnicos y científicos.
En la ingenieria mecatronica se usa frecuentemente para reportes,artículos y documentación de proyectos

\section{¿Por que usar \LaTeX?}
Algunas ventajas de usar \LaTeX son:
\begin{itemize}
 
    \item Excelente manejo de ecuaciones
    \item Documentos con formato profesional
    \item Ideal para trabajos de ingeniería 

\end{itemize}

\section{Ejemplo de ecuacion}
Una ecuacion común en mecatronica es la segunda ley de newton:
\begin{equation}
    F= m \cdot a
\end{equation}
otra ecuacion utilizada es la ley de ohm, la cual relaciona el corriente el volatje y la resistencia en uncircuito electrico: 

\begin{equation}
    V=I \cdot R
\end{equation}

\begin{equation}    
    I = \frac{V}{R}
\end{equation}

donde $I$ es la corriente eléctrica, $V$ es elvoltaje y $R$ es la resistencia, los cuales se miden en: Amperios($A$),voltios($V$) y Ohms($\Omega$) 

\section{Uso de imagenes}
Para importar imagenes en \LaTeX el comando a utilizar tiene una estructura muy sencilla 

\begin{figure}
\centering
\includegraphics[width=0.5\linewidth]{i1.jpg}
\caption{IMAGEN 1}
\label{fig:placeholder}
\end{figure}

\section{Conclusión}
\LaTeX es una herramienta poderosa que nos ayudara a presentar nuestros trabajos de forma clara y profesional.

\end{document}


